\begin{table}[htbp]
\centering
\caption{Marco macroeconómico, 2024--2025}
\label{tab:C1_1}
\small
\setlength{\tabcolsep}{4pt}
\begin{tabular}{L{6.0cm}ccc}
\toprule
{Variable} & {2024 aprobado} & {2024 estimado} & {2025} \\
\midrule
{PIB, crecimiento real (rango)} & {[2.5,3.5]} & {[1.5,2.5]} & {[2.0,3.0]} \\
{PIB nominal (miles de millones de pesos)} & \num{34374} & \num{33927.7} & \num{36166.4} \\
{Deflactor del PIB (\%)} & \num{4.8} & \num{4.6} & \num{4.3} \\
{Inflación dic/dic (\%)} & \num{3.8} & \num{4.3} & \num{3.5} \\
{Inflación promedio (\%)} & \num{4.5} & \num{4.7} & \num{3.8} \\
{Tipo de cambio fin de periodo} & \num{17.6} & \num{19.7} & \num{18.5} \\
{Tipo de cambio promedio} & \num{17.1} & \num{18.2} & \num{18.7} \\
{Cetes 28 nominal fin de periodo (\%)} & \num{9.5} & \num{10.0} & \num{8.0} \\
{Cetes 28 nominal promedio (\%)} & \num{10.3} & \num{10.7} & \num{8.9} \\
{Cetes 28 real acumulada (\%)} & \num{6.7} & \num{6.7} & \num{5.6} \\
{Cuenta corriente (millones de dólares)} & \num{-14954} & \num{-7097.5} & \num{-7941.0} \\
{Cuenta corriente (\% del PIB)} & \num{-0.7} & \num{-0.4} & \num{-0.4} \\
{PIB EE.UU., crecimiento real (\%)} & \num{1.8} & \num{2.7} & \num{2.2} \\
{Producción industrial EE.UU. (\%)} & \num{2.0} & \num{0.5} & \num{2.0} \\
{Inflación EE.UU. promedio (\%)} & \num{2.4} & \num{2.9} & \num{2.2} \\
{SOFR 3 meses promedio (\%)} & \num{4.3} & \num{5.0} & \num{3.3} \\
{Fed Funds Rate promedio (\%)} & \num{4.5} & \num{5.3} & \num{3.8} \\
{Petróleo, precio promedio (dls/barril)} & \num{56.7} & \num{70.7} & \num{57.8} \\
{Plataforma de producción (mbd)} & \num{1983} & \num{1877} & \num{1891} \\
{Plataforma de exportación (mbd)} & \num{994} & \num{900} & \num{892} \\
{Plataforma de privados (mbd)} & \num{86} & \num{67} & \num{74} \\
{Gas, precio promedio (dls/MMBtu)} & \num{3.5} & \num{2.2} & \num{3.1} \\
\bottomrule
\end{tabular}
\par\vspace{2pt}\footnotesize Fuente: elaboración propia del ITED con información de la SHCP, CGPE 2025, Anexo II.5, p. 81. El CGPE 2025 se publicó sin capa de texto; las cifras se leyeron de la página renderizada y se validaron por identidad contable.
\end{table}
